\section{The signed momentum estimate and its consequence}\label{sec:setup}
Our analytic goal is an estimate on momentum increments along
characteristics. We state it first and explain why its logarithmic time
scale excludes finite-time breakdown. Its proof occupies the central
sections of the paper.

Fix a finite horizon $T_0>0$ and a compact interval
$[0,t_*]\subset[0,T_{\max})$, with $t_*\le T_0$, of a classical solution.
A characteristic with initial label $z=(y_{\rm in},v_{\rm in})$ solves
\[
 Y'(s,z)=u(V(s,z)),\qquad
 V'(s,z)=E(s,Y(s,z))+u(V(s,z))\times B(s,Y(s,z)).
\]
It is \emph{supported} if $z\in\supp f_0$.
A distinguished supported characteristic, along which we estimate the
force, is written $(X(t),V_X(t))$ and called the \emph{receiver}.
We write
\[
 a(t)=u(V_X(t)),\quad q_X(t)=q(V_X(t)),\quad K(t)=V_X'(t).
\]
The signed increment is the vector impulse
\[
 V_X(t_2)-V_X(t_1)=\int_{t_1}^{t_2}K(t)\,dt.
\]
Source characteristics continue to use the letters $Y,V,u$.

For nonempty $\supp f_0$, choose a dyadic number $P$ such that
\begin{equation}\label{eq:P-range}
 P\ge1024\max_{\substack{0\le t\le t_*\\z\in\supp f_0}}q(V(t,z)),
 \qquad L=\log(2+P),\qquad S(w)=\frac{P^2L}{w}.
\end{equation}
A dyadic number is a power of two. Receiver scales are $w=2^k$, $k\ge0$.
In every nonempty range $w/8\le q_X\le8w$ one has $w\le P$.

\begin{proposition}[Signed momentum increments]\label{prop:increment}
For each smooth admissible datum and finite $T_0$, there are constants
$M,A\ge1$ and $P_0$ with the following property. On every compact
classical solution interval $[0,t_*]$ as above, for every dyadic
$P\ge P_0$ satisfying \eqref{eq:P-range}, every supported receiver,
every $w=2^k$, $k\ge0$, and every interval $[t_1,t_2]\subset[0,t_*]$
on which $w/8\le q_X(t)\le8w$ throughout, one has
\begin{equation}\label{eq:bootstrap}
 |V_X(t_2)-V_X(t_1)|
 \le MP\sqrt{t_2-t_1}+A(t_2-t_1)S(w).
\end{equation}
The constants are independent of $t_*$ and $P$.
\end{proposition}

\begin{proposition}[Local theory and continuation]\label{prop:continuation}
Every smooth admissible datum has a unique maximal classical solution.
It is smooth on each compact subinterval of its lifespan. If its
maximal time $T_{\max}$ is finite and
\[
 \sup\{\,|v|:0\le t<T_{\max},\ (x,v)\in\supp f(t)\,\}<\infty,
\]
then the solution has a classical extension beyond $T_{\max}$.
\end{proposition}

The proof is given in the final section. It combines the Glassey--Strauss
bounded-momentum criterion in the $H^5$ formulation quoted by Luk and
Strain~\cite[Theorem~1.1 and Footnote~1, arXiv version]{LukStrain2014Continuation}
with a finite-horizon modification of the initial fields. The modification retains the initial Coulomb field
and cuts off vector potentials for the divergence-free remainders.
The removed part evolves as a vacuum Maxwell field that vanishes
throughout the particle region for the chosen horizon. This supplies
Sobolev field data for the continuation theorem while preserving
the original particle dynamics.

\begin{proof}[Proof of Theorem~\ref{thm:global} from
Propositions~\ref{prop:increment} and~\ref{prop:continuation}]
If $f_0=0$, the vacuum Maxwell solution is global and smooth, as shown
in the final section. Suppose $\supp f_0\ne\varnothing$ and, toward a
contradiction, $T_{\max}<\infty$. Proposition~\ref{prop:continuation}
then forces unbounded momentum support. Set $T_0=T_{\max}$ and
\[
 Q(t)=\max_{z\in\supp f_0}q(V(t,z)).
\]
Continuity of the characteristic flow and compactness of the label
set make $Q$ continuous. For all sufficiently large integers $n$,
let $t_n<T_{\max}$ be the first time when $Q(t_n)=2^n$.
Choose a label attaining this value. Along that label let $s_n$
be its last time before $t_n$ at energy $2^{n-1}$. Such a time
exists because the initial energy is smaller; moreover
$s_n\ge t_{n-1}$. On $[s_n,t_n]$ the energy lies in
$[2^{n-1},2^n]$.

Apply Proposition~\ref{prop:increment} on the prefix $[0,t_n]$
with $P=1024\cdot2^n$ and $w=2^n$. Since $q$ is $1$-Lipschitz,
writing $\Delta_n=t_n-s_n$ gives
\[
 2^{n-1}\le |V(t_n)-V(s_n)|
 \le 1024M\,2^n\sqrt{\Delta_n}
   +1024^2A\,2^n\Delta_n\log(2+1024\cdot2^n).
\]
At least one term on the right is at least $2^{n-2}$.
Consequently
\[
 t_n-t_{n-1}\ge\Delta_n
 \ge \frac{c}{\log(2+1024\cdot2^n)}
\]
for a fixed $c>0$. The sum of these lower bounds diverges,
contradicting $t_n<T_{\max}$. Thus $T_{\max}=\infty$.
Uniqueness, smoothness and propagation of the constraints follow
from Proposition~\ref{prop:continuation} and its proof.
\end{proof}

To prove Proposition~\ref{prop:increment}, we assume \eqref{eq:bootstrap} for all such
receivers, scales and intervals in $[0,\tau]$, where $\tau\le t_*$,
and derive a uniform strict improvement. Until that improvement is
proved, \eqref{eq:bootstrap} is a bootstrap assumption.
In the estimates, $C$ and the constants implicit in $\lesssim$ or
$\asymp$ may depend on the datum, $T_0$, and fixed decomposition
parameters, but not on $t_*,P,M,A$. Additional dependence is indicated
explicitly. Lower thresholds on $P$ may depend on $M,A$.
The final choices are made in the order: a time-ramp fraction,
then $M$, then $A$, then the threshold on $P$.

For nonnegative estimates we allow integration over any measurable
subset of receiver times in an interval $J\subset[0,\tau]$ of length
at most $I$, provided those times satisfy $w/8\le q_X\le8w$ and
\begin{equation}\label{eq:short-I}
 0<I\le (w/P)^2.
\end{equation}
The number $I$ is an upper bound on length in these estimates.
For the final signed estimate the interval has length exactly $I$
and the receiver stays in the indicated range throughout.
This distinction permits summation over disjoint time pieces without
introducing a new constant for each piece.

Figure~\ref{fig:dependencies} records the conditional dependencies.
We now prove Proposition~\ref{prop:increment}. Energy and timelike
geometry give preliminary bounds on absolute force. Those bounds
lose a factor at large receiver momentum; the subsequent signed
cancellation removes this loss.

\begin{figure}[htbp]
\centering
\begin{tikzpicture}[
  x=1cm,y=1cm,
  stage/.style={draw=black!65,rounded corners=2pt,fill=white,
    line width=.55pt,align=center,text width=4.6cm,
    minimum height=.85cm,inner sep=5pt,font=\small},
  flow/.style={-{Stealth[length=2mm]},line width=.65pt,draw=black!75},
  note/.style={font=\footnotesize,fill=white,inner sep=2pt}
]
\node[stage,dashed] (assume) at (0,0)
  {Assumed increment bound\\up to time $\tau$};
\node[stage] (identity) at (7,0)
  {Signed cancellation\\along source labels};
\node[stage] (baseline) at (0,-1.5)
  {Baseline occupation\\and direct estimates};
\node[stage] (direction) at (0,-3)
  {Bound on changes\\of particle direction};
\node[stage] (occupation) at (0,-4.5)
  {Improved occupation\\in the angular window};
\node[stage] (selected) at (7,-4.5)
  {Selected-bin estimate\\for the signed force};
\node[stage] (strict) at (7,-6)
  {Strict improvement\\of the assumed bound};
\node[stage] (close) at (0,-6)
  {Continuity argument\\establishes the bound};
\draw[flow,dashed] (assume) -- (baseline);
\draw[flow] (baseline) -- (direction);
\draw[flow] (identity.south west) --
  node[note,pos=.44,above,sloped] {first use} (direction.east);
\draw[flow] (identity) --
  node[note,pos=.52,right] {second use} (selected);
\draw[flow] (direction) -- (occupation);
\draw[flow] (occupation) -- (selected);
\draw[flow] (selected) -- (strict);
\draw[flow] (strict) -- (close);
\end{tikzpicture}

\caption{Main dependencies in the conditional proof of the signed momentum estimate.
The dashed box is a bootstrap assumption. The first use of signed
cancellation controls particle direction from this assumption and
the baseline estimates; improved occupation is obtained afterward.
The second use yields a strict improvement, and a continuity argument
establishes the initially assumed bound.}
\label{fig:dependencies}

\end{figure}
